\documentclass[11pt,a4paper]{article}

% ============================================================================
% Packages
% ============================================================================
\usepackage[utf8]{inputenc}
\usepackage[T1]{fontenc}
\usepackage[italian,english]{babel}
\usepackage{geometry}
\geometry{margin=2.2cm}
\usepackage{amsmath,amssymb,amsfonts}
\usepackage{graphicx}
\usepackage{xcolor}
\usepackage{listings}
\usepackage{hyperref}
\usepackage{enumitem}
\usepackage{booktabs}
\usepackage{fancyhdr}
\usepackage{tcolorbox}
\usepackage{tikz}
\usetikzlibrary{arrows.meta,shapes,positioning}

% ============================================================================
% Formatting
% ============================================================================
\hypersetup{colorlinks=true,linkcolor=blue!70!black,urlcolor=blue!60!black,citecolor=green!50!black}

\definecolor{codegreen}{rgb}{0,0.5,0}
\definecolor{codegray}{rgb}{0.5,0.5,0.5}
\definecolor{codebg}{rgb}{0.97,0.97,0.97}
\definecolor{warnbg}{rgb}{1.0,0.95,0.85}
\definecolor{tipbg}{rgb}{0.88,0.97,0.88}
\definecolor{theorybg}{rgb}{0.90,0.92,1.0}
\definecolor{resultbg}{rgb}{1.0,0.96,0.96}

\lstset{
  language=C++,
  basicstyle=\ttfamily\small,
  backgroundcolor=\color{codebg},
  keywordstyle=\color{blue!70!black}\bfseries,
  commentstyle=\color{codegreen},
  stringstyle=\color{red!60!black},
  numbers=left, numberstyle=\tiny\color{codegray},
  frame=single, framerule=0.4pt,
  breaklines=true,
  tabsize=4,
  showstringspaces=false,
  columns=flexible
}

\newtcolorbox{warnbox}[1][]{colback=warnbg,colframe=orange!80!black,title={\textbf{Attenzione}},#1}
\newtcolorbox{tipbox}[1][]{colback=tipbg,colframe=green!60!black,title={\textbf{Suggerimento}},#1}
\newtcolorbox{theorybox}[1][]{colback=theorybg,colframe=blue!60!black,title={\textbf{Teoria}},#1}
\newtcolorbox{resultbox}[1][]{colback=resultbg,colframe=red!60!black,title={\textbf{Risultato misurato}},#1}

\pagestyle{fancy}
\fancyhf{}
\lhead{SPM Modulo 1 --- Guida Implementativa}
\rhead{A.Y. 2025/26}
\cfoot{\thepage}

% ============================================================================
\title{\textbf{Guida Implementativa Completa\\SPM Modulo 1: Partition Mapping Kernel}\\[0.5em]
\large Vectorizzazione del Kernel di Mapping Chiave$\to$Partizione}
\author{Corso SPM --- A.Y. 2025/26}
\date{Ultimo aggiornamento: \today}
% ============================================================================

\begin{document}
\maketitle
\tableofcontents
\newpage

% ############################################################################
\section{Panoramica del Progetto}
% ############################################################################

\subsection{Obiettivo}
Implementare un kernel ad alte prestazioni che mappa un array di $N$ chiavi \texttt{uint64\_t} ad un array di partition identifier $\texttt{part\_id}[i] \in [0, P)$, dove $P$ è una potenza di 2. Questo è il primo stadio di un \textit{partitioned hash join with duplicates}.

\textbf{Citazione dalla traccia}: ``\textit{The goal is not only to obtain speedups, but also to understand what limits performance.}''

\subsection{Deliverables richiesti dalla traccia}
\begin{enumerate}
    \item \textbf{Plain C++} (senza intrinsics): due binari dallo stesso sorgente:
    \begin{itemize}
        \item \textbf{baseline}: compilato con auto-vectorization \textbf{disabilitata}
        \item \textbf{autovec}: compilato con auto-vectorization \textbf{abilitata}
        \item Evidenza della vettorizzazione tramite report GCC
    \end{itemize}
    \item \textbf{AVX2 intrinsics}: produce output identico alla versione plain
    \item \textbf{(Opzionale) CUDA}: kernel-only time + transfer time separati; verifica correttezza vs CPU
    \item \textbf{Report PDF} (max 4 pagine): scelte progettuali, evidence auto-vec, strategia intrinsics, risultati, analisi bottleneck
\end{enumerate}

\subsection{Struttura del progetto}
\begin{verbatim}
module_1/
├── Makefile                  # Build system (arch-detection automatica)
├── README.md                 # Istruzioni complete
├── include/common.hpp        # tipi, hash, keygen, timing, checksum, mmap
├── src/
│   ├── plain.cpp             # C++ plain → baseline + autovec
│   ├── avx2.cpp              # AVX2 intrinsics + verifica integrata
│   ├── cuda_kernel.cu        # CUDA con timing separati
│   └── dataset_creator.cpp   # Generatore dataset binari
├── scripts/
│   ├── run_bench.sh          # SLURM batch CPU (3 esperimenti)
│   └── run_cuda.sh           # SLURM batch CUDA
├── analysis/
│   ├── parse_results.py      # Parser output → CSV
│   └── plot_results.py       # Generatore 16 grafici
├── results/                  # CSV, grafici, vec reports
├── report/                   # PDF report finale (3 pagine)
└── guide/                    # Questa guida
\end{verbatim}


% ############################################################################
\section{Fondamenti Teorici e Riferimenti alle Lezioni}
\label{sec:teoria}
% ############################################################################

\subsection{Il collo di bottiglia di Von Neumann (L5\&6, slide 3--5)}

\begin{theorybox}[title=\textbf{Von Neumann Bottleneck}]
In un'architettura a memoria condivisa, il tempo di esecuzione di un kernel è limitato dal \textbf{più lento} tra:
\begin{itemize}
    \item $t_\text{comp}$: tempo di calcolo (operazioni aritmetiche/logiche)
    \item $t_\text{mem}$: tempo di trasferimento dati (memoria $\leftrightarrow$ processore)
\end{itemize}
Se $t_\text{mem} > t_\text{comp}$, il kernel è \textbf{memory-bound}: il processore è in stallo aspettando i dati dalla memoria. In questo caso, aumentare la potenza di calcolo (es.\ SIMD più largo) \textbf{non} migliora le prestazioni.
\end{theorybox}

\textbf{Il nostro kernel è memory-bound.} Per ogni chiave:
\begin{itemize}
    \item \textbf{Legge} 8 byte (una chiave \texttt{uint64\_t})
    \item \textbf{Scrive} 4 byte (un \texttt{uint32\_t} partition ID)
    \item \textbf{Calcola} $\sim$4 operazioni intere (XOR, mul32, shift, store)
\end{itemize}
Il traffico di memoria totale è $12N$ byte, il compute è $\sim 4N$ operazioni.


\subsection{Modello Roofline (L5\&6, slide 47--51)}

\begin{theorybox}[title=\textbf{Modello Roofline}]
La performance massima raggiungibile (in ops/s) è:
$$P_\text{max} = \min\bigl(R_\text{peak},\; I \times B_\text{peak}\bigr)$$
dove:
\begin{itemize}
    \item $R_\text{peak}$: prestazione computazionale di picco del processore (ops/s)
    \item $B_\text{peak}$: bandwidth di memoria di picco (byte/s)
    \item $I = \frac{\text{operazioni}}{\text{byte trasferiti}}$: \textbf{Operational Intensity} (ops/byte)
\end{itemize}
Il \textbf{ridge point} $I^* = R_\text{peak} / B_\text{peak}$ separa la regione memory-bound ($I < I^*$) dalla regione compute-bound ($I > I^*$).
\end{theorybox}

\paragraph{Calcolo per il nostro kernel.}
\begin{align}
I &= \frac{4 \text{ ops}}{12 \text{ byte}} = 0.333 \text{ ops/byte} \label{eq:oi}
\end{align}

\begin{resultbox}[title=\textbf{Parametri hardware misurati su node09 (micro-benchmark)}]
\begin{itemize}
    \item $B_\text{peak} = 16.4$ GB/s (single-core, misurato con STREAM-like AVX2 copy)
    \item $R_\text{peak}^\text{scalar} = 5.17$ Gops/s (4 catene indipendenti XOR+MUL+SHIFT)
    \item $R_\text{peak}^\text{AVX2} = 15.28$ Gops/s (4 catene $\times$ 8 elementi)
\end{itemize}
Ridge point scalare: $I^* = 5.17 / 16.4 = 0.315$ ops/byte.\\
Ridge point AVX2: $I^* = 15.28 / 16.4 = 0.932$ ops/byte.

Il nostro kernel ha $I = 0.333$, che è:
\begin{itemize}
    \item Appena \textbf{sopra} il ridge point scalare ($0.333 > 0.315$) → leggermente compute-bound in scalare
    \item \textbf{Sotto} il ridge point AVX2 ($0.333 < 0.932$) → fortemente memory-bound con SIMD
\end{itemize}
Questo spiega perché il passaggio da scalare a AVX2 sposta il bottleneck da compute a memoria.
\end{resultbox}

Il tetto di bandwidth al nostro $I$:
$$P_\text{bw} = I \times B_\text{peak} = 0.333 \times 16.4 = 5.47 \text{ Gops/s}$$

L'autovec misurato raggiunge 5.24 Gops/s = \textbf{96\% del tetto di bandwidth}. Il kernel è essenzialmente \textit{bandwidth-saturated}.


\subsection{Modello SIMD (L7\&8, slide 7)}

\begin{theorybox}[title=\textbf{Speedup SIMD teorico}]
$$\text{Speedup}_\text{SIMD} \approx \text{vector\_width} \times \text{efficiency}$$
Con AVX2 su uint32: vector\_width = 8, ma l'efficienza è limitata dalla memoria.
\end{theorybox}

\paragraph{L7\&8, slide 8}: ``\textit{Not all intrinsic functions map one-to-one to a single instruction}''. Questo è cruciale per la nostra scelta di hash: una \texttt{\_mm256\_mullo\_epi64} non esiste in AVX2 (solo AVX-512). La decomposizione della mul64 richiede 3$\times$ \texttt{vpmuludq} + shift + add $\Rightarrow$ la SIMD diventa più lenta dello scalare.

\paragraph{L7\&8, slide 30, 39}: SoA layout e stride unitario sono prerequisiti per la vettorizzazione efficiente. Il nostro kernel accede a \texttt{keys[i]} e \texttt{part\_ids[i]} in modo sequenziale.

\paragraph{L7\&8, slide 32--34, 40}: condizioni per l'auto-vectorization di GCC.


\subsection{Metriche di Performance (L10)}

\begin{theorybox}[title=\textbf{Definizioni formali (L10, slide 5--6)}]
\begin{itemize}
    \item \textbf{Speedup}: $S(p) = T_\text{seq} / T_\text{par}(p)$
    \item \textbf{Efficienza}: $E(p) = S(p) / p$
    \item \textbf{Costo}: $C(p) = T_\text{par}(p) \times p$
\end{itemize}
Nel contesto SIMD, $p$ è la \textbf{larghezza del vettore}: $p=1$ (scalare), $p=8$ (AVX2 256-bit su uint32).
\end{theorybox}

\paragraph{Legge di Amdahl (L10, slide 24--28).}
$$S(p) \leq \frac{1}{f + (1-f)/p}$$
dove $f$ è la \textbf{frazione seriale} del programma. Per $p \to \infty$: $S_\text{max} = 1/f$.

\textbf{Interpretazione per SIMD memory-bound}: la ``frazione seriale'' $f$ non è codice non-parallelizzabile, ma il \textbf{tempo di attesa dalla memoria}, che non beneficia del parallelismo SIMD. Dal nostro speedup misurato $S = 1.46\times$ con $p=8$:
$$f = \frac{1/S - 1/p}{1 - 1/p} = \frac{1/1.46 - 1/8}{1 - 1/8} \approx 64\%$$
Anche con $p \to \infty$: $S_\text{max} \approx 1/0.64 = 1.56\times$. Passare ad AVX-512 ($p=16$) non darebbe guadagni significativi.

\paragraph{Legge di Amdahl con overhead (L10, slide 31--32).}
La legge di Amdahl ``pura'' è ottimistica perché ignora gli overhead. Nel nostro caso SIMD:
$$S(p) \leq \frac{1}{[f + O(p)] + (1-f)/p}$$
dove $O(p)$ include: costo delle istruzioni di permutazione/pack, tail loop per $N \bmod 8$ chiavi.


\subsection{CUDA e modello SIMT (L9)}

\begin{theorybox}[title=\textbf{Modello SIMT (L9, slide 2--5)}]
GPU NVIDIA eseguono migliaia di thread raggruppati in warp da 32 thread. Ogni thread esegue la stessa istruzione. Per kernel memory-bound, la chiave è la \textbf{bandwidth HBM2} e il \textbf{memory coalescing} (L9, slide 28).
\end{theorybox}

\paragraph{A30 specs (L9, slide 24)}: 993 GB/s HBM2 peak, PCIe 4.0 x16 ($\sim$25 GB/s bidirezionale).

\paragraph{Il problema del trasferimento}: per un kernel che fa pochissimo compute per byte, il trasferimento PCIe Host$\leftrightarrow$Device domina il tempo end-to-end. Il kernel GPU è utile solo se i dati sono \textbf{già residenti in GPU memory}.


\subsection{Hashing: dal problema alla scelta}

\begin{theorybox}[title=\textbf{Il problema dell'hashing per il partitioning}]
Data una chiave $k \in U = \{0, \ldots, 2^{64}-1\}$ e un numero di partizioni $P$, vogliamo $h: U \to \{0, \ldots, P-1\}$ che sia:
\begin{itemize}
    \item \textbf{Deterministica}: stessa chiave $\to$ stessa partizione
    \item \textbf{Uniformemente distribuita}: chiavi equamente tra le $P$ partizioni
    \item \textbf{Veloce}: applicata a \emph{tutti} gli $N$ input (streaming kernel)
    \item \textbf{SIMD-friendly}: operazioni vettorizzabili nativamente
\end{itemize}
\end{theorybox}

\subsubsection{Hashing by Division ($k \bmod P$)}
Il metodo più semplice. \textbf{Problema}: l'operazione modulo è costosa (divisione intera), e se $P$ è potenza di 2 si riduce a una maschera di bit ($k\ \&\ (P-1)$), ma questo considera solo i bit \textbf{meno significativi} della chiave --- distribuzione pessima per chiavi con pattern nei bit bassi.

\subsubsection{Universal Hashing e Multiply-Add-Shift}
Dalla classe universale Multiply-Add-Shift (Dietzfelbinger et al., 1997):
$$h_{a,b}(k) = \left\lfloor \frac{(ak + b) \bmod 2^w}{2^{w-\ell}} \right\rfloor$$
dove $a$ è un intero \textbf{dispari} $< 2^w$ e $b \geq 0$. Questa classe è \textbf{2-universale}: per ogni coppia di chiavi distinte $x \neq y$, la probabilità di collisione è al più $1/m$.

\paragraph{Semplificazione $b=0$ (Multiply-Shift).}
Ponendo $b = 0$ si ottiene la variante \textit{quasi-universale} (collisione $\leq 2/m$):
$$h_a(k) = (A \cdot k) \gg (w - \log_2 P)$$
Per il partizionamento (non una hash table con chaining), la garanzia $\leq 2/P$ è ampiamente sufficiente. In cambio, risparmiamo un'addizione per chiave --- significativo in un kernel streaming su centinaia di milioni di chiavi.

\paragraph{La costante $A$: perché il golden ratio.}
La costante $A$ non basta che sia dispari --- è necessario (per garantire biiezione in $\mathbb{Z}_{2^w}$), ma non sufficiente per buona distribuzione pratica ($a=1$ e $a=3$ sono dispari ma pessimi). La scelta $A = \lfloor 2^w / \varphi \rfloor$ (dove $\varphi = (1+\sqrt{5})/2$) ha tre proprietà speciali:
\begin{enumerate}
    \item \textbf{Equidistribuzione di Weyl}: la sequenza $\{k \cdot \alpha\}$ (parte frazionaria) è equidistribuita mod~1 per ogni irrazionale $\alpha$. Il reciproco del golden ratio $1/\varphi$ è il numero ``più irrazionale'' --- convergenza più lenta con frazioni continue $\Rightarrow$ massima uniformità (Knuth, TAOCP Vol.~3, §6.4).
    \item \textbf{Bit-mixing}: la rappresentazione binaria di \texttt{0x9E3779B9} ha una miscela quasi bilanciata di 0 e 1 distribuiti irregolarmente. Ogni bit del risultato dipende da molti bit dell'input.
    \item \textbf{Dispari per costruzione}: $\lfloor 2^{32}/\varphi \rfloor$ è dispari, soddisfacendo automaticamente il requisito formale.
\end{enumerate}

\subsubsection{La nostra scelta: XOR-fold + Fibonacci mul32}

\begin{equation}
\boxed{h(k) = \bigl((k_\text{lo} \oplus k_\text{hi}) \cdot A_{32}\bigr) \gg (32 - \log_2 P)}
\label{eq:hash}
\end{equation}
dove $A_{32} = \texttt{0x9E3779B9} = \lfloor 2^{32}/\varphi \rfloor$ (costante Fibonacci, Knuth TAOCP Vol.~3, §6.4).

\textbf{Motivazioni}:
\begin{enumerate}
    \item \textbf{SIMD-native}: \texttt{\_mm256\_mullo\_epi32} (\texttt{vpmulld}) è nativa AVX2 --- 8 mul/istruzione
    \item \textbf{Entropy preservation}: $k_\text{lo} \oplus k_\text{hi}$ combina entrambe le metà della chiave
    \item \textbf{Golden ratio}: la costante $\lfloor 2^{32}/\varphi \rfloor$ ha proprietà speciali di equidistribuzione (Weyl) e bit-mixing ottimale
    \item \textbf{Minimal compute}: solo XOR + mul32 + shift (no divisione, no modulo, no branch)
    \item \textbf{$P$ potenza di 2}: lo shift sostituisce il modulo
\end{enumerate}

\begin{warnbox}[title=\textbf{Perché non una hash a 64 bit?}]
La scelta naturale sarebbe $h(k) = (A_{64} \cdot k) \gg (64 - \log_2 P)$ con un'istruzione scalare \texttt{imul r64,r64}. Ma \textbf{AVX2 non ha \texttt{mullo\_epi64}} (disponibile solo da AVX-512). La decomposizione richiede 3$\times$ \texttt{vpmuludq} + 2 shift + 2 add per 4 chiavi.

Come nota la slide 8 di L7\&8: ``\textit{not all intrinsic functions map one-to-one to a single instruction}''.

\textbf{Risultato misurato su node09}: con hash mul64, AVX2 è \textbf{più lento} dello scalare ($\sim$880 vs $\sim$910 Mkeys/s). Con XOR-fold + mul32, AVX2 raggiunge $\sim$1200 Mkeys/s (1.31$\times$ speedup).
\end{warnbox}


% ############################################################################
\section{Implementazione Dettagliata}
% ############################################################################

\subsection{common.hpp --- Infrastruttura condivisa}

Il file \texttt{include/common.hpp} contiene tutti i componenti riusati:

\paragraph{Tipi.}
\begin{lstlisting}
using spm_key_t = uint64_t;  // chiave 64-bit
using part_t    = uint32_t;  // partition ID
\end{lstlisting}
\textbf{Nota}: usiamo \texttt{spm\_key\_t} invece di \texttt{key\_t} per evitare il conflitto con la \texttt{key\_t} POSIX definita in \texttt{sys/types.h}.

\paragraph{Funzione hash.}
\begin{lstlisting}
static constexpr uint32_t HASH_A32 = 0x9E3779B9u;

inline part_t hash_key(spm_key_t key, unsigned shift32) {
    uint32_t k_lo = static_cast<uint32_t>(key);
    uint32_t k_hi = static_cast<uint32_t>(key >> 32);
    return (uint32_t)(((k_lo ^ k_hi) * HASH_A32) >> shift32);
}
\end{lstlisting}

\paragraph{Generatore chiavi deterministico.} Usa xoshiro256** (Blackman--Vigna), inizializzato con SplitMix64 dal seed. Il parametro \texttt{key\_space} permette di controllare i duplicati ($\texttt{key} \bmod \texttt{key\_space}$).

\paragraph{Checksum FNV-1a: perché funziona come verifica di correttezza.}

La traccia richiede ``\textit{a verification strategy that can detect mismatches without printing the full output array}''. Usiamo FNV-1a (Fowler--Noll--Vo), un hash non crittografico progettato per avere ottime proprietà di \textit{avalanche}:

\begin{lstlisting}
inline uint64_t compute_checksum(const part_t* part_ids, size_t N) {
    uint64_t h = 0xCBF29CE484222325ULL;  // offset basis
    for (size_t i = 0; i < N; i++) {
        h ^= static_cast<uint64_t>(part_ids[i]);
        h *= 0x100000001B3ULL;  // FNV prime
    }
    return h;
}
\end{lstlisting}

\begin{theorybox}[title=\textbf{Perché il checksum assicura correttezza}]
\begin{enumerate}
    \item \textbf{Sensibilità alla posizione}: FNV-1a è \textit{order-dependent} --- XOR + moltiplicazione ad ogni passo fanno sì che cambiare anche un solo \texttt{part\_ids[i]} o scambiarne due produca un hash completamente diverso (effetto valanga). Non è un semplice XOR o somma degli elementi (che sarebbero commutative e insensibili all'ordine).
    \item \textbf{Probabilità di falso positivo}: con un output a 64 bit, la probabilità che due array diversi producano lo stesso checksum è $\lesssim 2^{-64} \approx 5.4 \times 10^{-20}$. In pratica: se i checksum coincidono, gli array sono identici.
    \item \textbf{Costo $O(N)$}: il checksum si calcola in un singolo passaggio sull'array, senza allocazione aggiuntiva. Per $N=10^8$ chiavi, sono $\sim$400 MB di dati --- stampare e confrontare sarebbe impraticabile.
    \item \textbf{Confronto cross-implementazione}: tutte le versioni (baseline, autovec, AVX2, CUDA) calcolano lo stesso checksum sullo stesso output. Se i 4 checksum coincidono, le 4 implementazioni producono \textbf{esattamente} lo stesso array.
\end{enumerate}
\end{theorybox}

\begin{warnbox}[title=\textbf{Limiti del checksum}]
Un checksum \textbf{non dimostra} che l'output sia ``corretto'' in senso assoluto --- dimostra solo che tutte le implementazioni producono lo \textbf{stesso} output. Se tutte avessero lo stesso bug, il checksum coinciderebbe comunque. Per questo la traccia richiede anche il confronto element-wise per $N$ piccolo e la verifica della distribuzione tra partizioni.
\end{warnbox}

\paragraph{Allocazione allineata.} Usa \texttt{posix\_memalign} (non \texttt{std::aligned\_alloc} che non è disponibile su macOS Apple libc++). L'allineamento a 32 byte è richiesto da \texttt{\_mm256\_load\_si256}.

\paragraph{Dataset binari e mmap.} I dataset sono file binari con header 40 byte (magic, N, seed, key\_space) + $N \times 8$ byte di chiavi. Caricati via \texttt{mmap()} per zero-copy.

\paragraph{Benchmark utilities.} Mediana di $r$ ripetizioni, deviazione standard, throughput in Mkeys/s.


\subsection{plain.cpp --- Kernel scalare (Task 1)}

\paragraph{Il kernel.}
\begin{lstlisting}
void partition_map(const spm_key_t* __restrict__ keys,
                   part_t*          __restrict__ part_ids,
                   size_t N, unsigned shift) {
    for (size_t i = 0; i < N; i++) {
        uint32_t k_lo = static_cast<uint32_t>(keys[i]);
        uint32_t k_hi = static_cast<uint32_t>(keys[i] >> 32);
        uint32_t mixed = k_lo ^ k_hi;
        part_ids[i] = (mixed * HASH_A32) >> shift;
    }
}
\end{lstlisting}

\textbf{Condizioni per auto-vectorization} (L7\&8, slide 32--34):
\begin{itemize}
    \item[\checkmark] Trip count noto ($N$ fisso all'ingresso del loop)
    \item[\checkmark] Nessuna dipendenza tra iterazioni (ogni \texttt{part\_ids[i]} dipende solo da \texttt{keys[i]})
    \item[\checkmark] Nessuna function call nel body (tutto inline)
    \item[\checkmark] \texttt{\_\_restrict\_\_} per escludere aliasing tra \texttt{keys} e \texttt{part\_ids}
    \item[\checkmark] Stride unitario per entrambi gli array
    \item[\checkmark] Operazione SIMD-native (\texttt{vpmulld} per mul32)
\end{itemize}

\paragraph{Due binari dallo stesso sorgente.}
\begin{lstlisting}[language=bash]
# Baseline (auto-vectorization disabilitata):
g++ -std=c++20 -O3 -march=native -mavx2 -fno-tree-vectorize \
    src/plain.cpp -I include -o bin/plain_baseline

# Autovec (auto-vectorization abilitata con report):
g++ -std=c++20 -O3 -march=native -mavx2 -ftree-vectorize \
    -DAUTOVEC_ENABLED \
    -fopt-info-vec-optimized=results/vec_report_optimized.txt \
    -fopt-info-vec-missed=results/vec_report_missed.txt \
    src/plain.cpp -I include -o bin/plain_autovec
\end{lstlisting}

\paragraph{Evidenza auto-vectorization (GCC 12.2 su node09).}
\begin{lstlisting}
src/plain.cpp:33:26: optimized: loop vectorized using 32 byte vectors
src/plain.cpp:33:26: optimized: loop vectorized using 16 byte vectors
\end{lstlisting}

GCC emette sia la versione a 256 bit (ymm) che a 128 bit (xmm). Su \textbf{Zen~1}, le operazioni AVX2 a 256 bit vengono decomposte in 2$\times$128 bit dalla microarchitettura (le ALU interne sono a 128 bit), ma \texttt{vpmulld} rimane un'istruzione nativa.

\paragraph{Il main.} Gestisce:
\begin{itemize}
    \item Parsing parametri: N, P, seed, key\_space, reps
    \item Validazione $P$ potenza di 2: \texttt{P \& (P-1) == 0}
    \item Generazione chiavi, warmup (1 esecuzione per popolare cache/TLB)
    \item Loop benchmark con \texttt{std::chrono::high\_resolution\_clock}
    \item Stampa: mediana, stddev, throughput, checksum
    \item Element-wise per $N \leq 32$
    \item Distribuzione tra partizioni per $P \leq 1024$
\end{itemize}


\subsection{avx2.cpp --- Intrinsics AVX2 (Task 2)}

\paragraph{Strategia: 8 chiavi per iterazione.}
Poiché le chiavi sono a 64 bit ma la hash opera a 32 bit, il flusso è:
\begin{enumerate}
    \item \textbf{Carica} 8 chiavi uint64 in 2 registri ymm ($4 \times 64 = 256$ bit ciascuno)
    \item \textbf{XOR-fold}: estrai k\_hi con \texttt{\_mm256\_srli\_epi64(vk, 32)}, XOR con k\_lo
    \item \textbf{Pack}: estrai i 32 bit bassi da ogni lane 64-bit, combina in un singolo ymm di $8 \times 32$ bit
    \item \textbf{Multiply}: \texttt{\_mm256\_mullo\_epi32} (una sola istruzione \texttt{vpmulld})
    \item \textbf{Shift}: \texttt{\_mm256\_srli\_epi32} per ottenere il partition ID
    \item \textbf{Store}: \texttt{\_mm256\_storeu\_si256} (8 uint32 = 256 bit)
\end{enumerate}

\paragraph{Istruzioni critiche per iterazione:}
\begin{itemize}
    \item 2 $\times$ \texttt{vmovdqa} (load 256-bit aligned)
    \item 2 $\times$ \texttt{vpsrlq} (shift-right 64-bit per estrarre k\_hi)
    \item 2 $\times$ \texttt{vpxor} (XOR fold)
    \item 2 $\times$ \texttt{vpermd} (permute per estrarre 32 bit bassi)
    \item 1 $\times$ \texttt{vperm2i128} (combina le due metà)
    \item 1 $\times$ \texttt{vpmulld} (Fibonacci mul32 --- 8 moltiplicazioni)
    \item 1 $\times$ \texttt{vpsrld} (shift per partition ID)
    \item 1 $\times$ \texttt{vmovdqu} (store 256-bit unaligned)
\end{itemize}
Totale: $\sim$12 istruzioni SIMD per 8 chiavi = \textbf{1.5 istruzioni/chiave}.\\
Scalare: $\sim$3 istruzioni per chiave (XOR + IMUL + SHR).\\
Rapporto istruzioni: $3/1.5 = 2\times$, ma il kernel è memory-bound (Sezione~\ref{sec:risultati}).

\paragraph{Tail loop.} Per gli ultimi $N \bmod 8$ elementi, usa la funzione \texttt{hash\_key()} scalare.

\paragraph{Verifica integrata.} Il binario \texttt{avx2} esegue internamente \textbf{sia} il kernel scalare (con \texttt{\#pragma GCC optimize("no-tree-vectorize")}) che quello AVX2, confronta i checksum FNV-1a, e riporta PASS/FAIL. In caso di mismatch, indica la prima posizione discordante.


\subsection{cuda\_kernel.cu --- CUDA (Task 3, opzionale)}

\paragraph{Kernel CUDA.}
\begin{lstlisting}
__global__
void partition_map_cuda(const uint64_t* __restrict__ keys,
                        uint32_t*       __restrict__ part_ids,
                        size_t N, uint32_t hash_a32, unsigned shift) {
    size_t idx = (size_t)blockIdx.x * blockDim.x + threadIdx.x;
    if (idx < N) {
        uint32_t k_lo = (uint32_t)(keys[idx]);
        uint32_t k_hi = (uint32_t)(keys[idx] >> 32);
        part_ids[idx] = ((k_lo ^ k_hi) * hash_a32) >> shift;
    }
}
\end{lstlisting}

Ogni thread mappa \textbf{una} chiave. Con 256 threads/block e accesso sequenziale a \texttt{keys[idx]}, i warp accedono a indirizzi consecutivi $\Rightarrow$ \textbf{memory coalescing ottimale} (L9, slide 28).

\paragraph{Timing separati.} La traccia richiede esplicitamente di riportare separatamente H$\to$D, kernel, D$\to$H. Usiamo \texttt{cudaEvent\_t} per misurazioni accurate:
\begin{lstlisting}
cudaEventRecord(e0);  // start
cudaMemcpy(d_keys, h_keys, ...HostToDevice);
cudaEventRecord(e1);  // dopo H->D
partition_map_cuda<<<nblocks, tpb>>>(...);
cudaEventRecord(e2);  // dopo kernel
cudaMemcpy(h_part, d_part, ...DeviceToHost);
cudaEventRecord(e3);  // dopo D->H
\end{lstlisting}

\paragraph{Pinned memory.} \texttt{cudaMallocHost} per i buffer host: abilita DMA asincrono e aumenta il throughput PCIe.

\paragraph{Compilazione.}
\begin{lstlisting}[language=bash]
# Sul nodo 9 (A30, compute capability 8.0):
nvcc -std=c++17 -O3 -I include \
     -gencode arch=compute_80,code=sm_80 \
     src/cuda_kernel.cu -o bin/cuda_kernel
\end{lstlisting}

\paragraph{Correttezza.} Il binario esegue la stessa hash su CPU come riferimento, confronta i checksum FNV-1a, e riporta PASS/FAIL con indicazione del primo mismatch.


% ############################################################################
\section{Ambiente di Esecuzione: Cluster spmcluster, node09}
\label{sec:cluster}
% ############################################################################

\subsection{Hardware di node09}
\begin{table}[h]
\centering
\begin{tabular}{@{}ll@{}}
\toprule
\textbf{Componente} & \textbf{Specifica} \\
\midrule
CPU & AMD EPYC 7301 (Zen 1 / Naples) \\
Sockets $\times$ Cores $\times$ Threads & 2 $\times$ 16 $\times$ 2 = 64 CPU \\
Boost clock & 2.7 GHz \\
ISA & x86\_64, AVX2, FMA (ALU 128-bit, AVX2 decomposto 2$\times$128) \\
RAM & 257 GB DDR4-2666, 8 nodi NUMA \\
Cache & L1d: 1 MiB, L2: 16 MiB, L3: 128 MiB \\
GPU & NVIDIA A30 (compute capability 8.0) \\
GPU memory & HBM2, 993 GB/s peak \\
Interconnect & PCIe 4.0 x16 ($\sim$25 GB/s bidirezionale) \\
\bottomrule
\end{tabular}
\end{table}

\begin{warnbox}[title=\textbf{Zen 1 e AVX2}]
L'EPYC 7301 (Zen 1) ha ALU interne a 128 bit. Le operazioni AVX2 a 256 bit vengono \textbf{decomposte in 2 micro-ops a 128 bit}. Questo significa che \texttt{vpmulld ymm} (8 $\times$ mul32) richiede 2 cicli (2 $\times$ 4 mul32), non 1. Tuttavia rimane \textbf{molto più efficiente} della decomposizione mul64 (che richiederebbe 3 $\times$ vpmuludq = 6 micro-ops per 4 chiavi).

Altre micro-architetture rilevanti:
\begin{itemize}
    \item \texttt{imul r64,r64}: latency 3 cicli, throughput 1/ciclo (Zen 1)
    \item \texttt{vpmulld} (128-bit): latency 4 cicli, throughput 1/ciclo
\end{itemize}
\end{warnbox}

\subsection{SLURM}
\begin{itemize}
    \item Partizioni verificate: \texttt{gpu-excl} (EXCLUSIVE, 30 min max), \texttt{gpu-shared} (10 min)
    \item I benchmark usano: \texttt{--partition=gpu-excl --nodelist=node09 --ntasks=1 --cpus-per-task=1}
    \item L'allocazione esclusiva elimina interferenze da altri job
\end{itemize}

\subsection{Compilatore}
\begin{itemize}
    \item GCC 12.2 (OpenHPC): compila i target CPU
    \item nvcc / CUDA 12.3 (in \texttt{/usr/local/cuda-12.3/}): compila CUDA \textbf{direttamente sul nodo 9} (lo script \texttt{run\_cuda.sh} include la compilazione perché nvcc è disponibile solo lì)
\end{itemize}


% ############################################################################
\section{Risultati Sperimentali}
\label{sec:risultati}
% ############################################################################

Tutti i benchmark su node09, allocazione esclusiva, 11 ripetizioni, mediana riportata.

\subsection{Tabella riassuntiva (N=100M, P=256)}

\begin{table}[h]
\centering
\begin{tabular}{@{}lcccc@{}}
\toprule
\textbf{Versione} & \textbf{Mediana (ms)} & \textbf{Mkeys/s} & \textbf{Speedup} & \textbf{BW (GB/s)} \\
\midrule
Baseline (no-vec) & 109.0 & 918 & 1.00$\times$ & 11.0 \\
Auto-vectorized   &  76.3 & 1310 & 1.43$\times$ & 15.7 \\
AVX2 intrinsics   &  83.3 & 1200 & 1.31$\times$ & 14.4 \\
\midrule
CUDA kernel-only  &   1.5 & 67\,320 & 73.3$\times$ & 808 \\
CUDA end-to-end   &  97.0 & 1\,031 & 1.12$\times$ & 12.4 \\
\bottomrule
\end{tabular}
\caption{BW = throughput $\times$ 12 B/key.}
\end{table}

\textbf{Correttezza}: tutti i checksum FNV-1a corrispondono tra baseline, autovec, AVX2 e CUDA per ogni combinazione $(N, P, \text{seed})$ testata (9 configurazioni verificate).

\subsection{Analisi Roofline}

\begin{resultbox}
\begin{itemize}
    \item Operational intensity: $I = 4/12 = 0.333$ ops/byte
    \item Tetto di bandwidth: $I \times B_\text{peak} = 0.333 \times 16.4 = 5.47$ Gops/s
    \item Autovec raggiunge 5.24 Gops/s = \textbf{96\% del tetto}
    \item AVX2 raggiunge 4.80 Gops/s = \textbf{88\% del tetto}
    \item Baseline raggiunge 3.67 Gops/s = \textbf{67\% del tetto}
\end{itemize}
Tutte le implementazioni sono nella \textbf{regione bandwidth-bound} del Roofline.
\end{resultbox}

\textbf{Perché lo speedup è solo 1.3--1.4$\times$ invece di 8$\times$?}
\begin{enumerate}
    \item AVX2 processa 8 uint32 per istruzione, ma questo accelera solo $t_\text{comp}$
    \item Il kernel è memory-bound ($t_\text{mem} > t_\text{comp}$ con SIMD)
    \item Lo speedup viene dal miglior utilizzo della bandwidth: meno istruzioni $\Rightarrow$ meno stalli frontend $\Rightarrow$ prefetch più efficiente
    \item Il tetto fisico è la bandwidth DRAM single-core (16.4 GB/s)
\end{enumerate}

\textbf{Perché autovec (1.43$\times$) $>$ AVX2 intrinsics (1.31$\times$)?}\\
GCC applica \textbf{loop unrolling} e \textbf{software pipelining} automatici che il codice intrinsics ``rigido'' non beneficia. Il compilatore ha più libertà di riordinare le istruzioni per nascondere la latenza di memoria.

\subsection{Metriche SIMD (L10)}

Applicando le metriche della lezione 10 con $p = 8$ (8 lane AVX2 su uint32):

\begin{table}[h]
\centering
\begin{tabular}{@{}lcc@{}}
\toprule
\textbf{Metrica} & \textbf{Autovec} & \textbf{AVX2} \\
\midrule
Speedup $S(8)$ & 1.46$\times$ & 1.27$\times$ \\
Efficienza $E(8) = S/p$ & 18.2\% & 15.9\% \\
Costo $C(8) = T_\text{par} \times p$ & $5.5 \times T_\text{seq}$ & $6.3 \times T_\text{seq}$ \\
$f$ (Amdahl) & 64.2\% & 75.8\% \\
$S_\text{max} = 1/f$ & 1.56$\times$ & 1.32$\times$ \\
\bottomrule
\end{tabular}
\caption{$N=200$M, $P=256$. Efficienza $\sim$16--18\%: le lane SIMD sono in stallo per la maggior parte del tempo, aspettando dati dalla memoria.}
\end{table}

Il \textbf{costo} $C(p) = T_\text{par} \times p$ è 5--6$\times$ superiore a $T_\text{seq}$: la parallelizzazione SIMD ``spreca'' risorse hardware perché la memoria non riesce ad alimentare tutte le lane contemporaneamente.

\subsection{Sweep dei parametri}

\paragraph{N-sweep (P=256).} Lo speedup è stabile da 1M a 200M chiavi. Questo conferma il comportamento bandwidth-bound: al crescere di $N$, baseline e versioni vettorizzate scalano linearmente, mantenendo il rapporto costante.

\paragraph{P-sweep (N=100M).} Throughput e speedup sono indipendenti da $P \in [2, 1024]$. Modificare $P$ cambia solo la costante di shift --- il traffico di memoria e il compute per chiave non cambiano.

\paragraph{Key-space sweep (N=100M, P=256).} Throughput invariante al tasso di duplicati (\texttt{key\_space} da $10^3$ a $10^9$ e 0=full range). La hash è branchless con costo costante per chiave.


\subsection{CUDA: breakdown dei tempi}

\begin{table}[h]
\centering
\begin{tabular}{@{}lrrc@{}}
\toprule
\textbf{Fase} & \textbf{Tempo (ms)} & \textbf{BW (GB/s)} & \textbf{\% totale} \\
\midrule
H$\to$D (800 MB) & 63.2 & 12.7 & 65.2\% \\
Kernel           &  1.5 & 808  &  1.5\% \\
D$\to$H (400 MB) & 32.3 & 12.4 & 33.3\% \\
\midrule
\textbf{Totale}  & \textbf{97.0} & --- & 100\% \\
\bottomrule
\end{tabular}
\caption{$N=100$M, $P=256$. Il kernel CUDA è il 1.5\% del tempo totale.}
\end{table}

\textbf{Kernel-only}: 67\,320 Mkeys/s = 808 GB/s = \textbf{81\% della peak bandwidth HBM2} (993 GB/s). Ottimo utilizzo della GPU.

\textbf{End-to-end}: $\sim$1\,031 Mkeys/s --- paragonabile alla CPU. Il trasferimento PCIe ($\sim$12.5 GB/s misurato, vs 25 GB/s teorico half-duplex) consuma il \textbf{98.5\% del tempo}.

\textbf{Conclusione}: GPU offload di questo kernel \emph{isolato} non conviene. Diventa vantaggioso solo se:
\begin{itemize}
    \item I dati sono già residenti in GPU memory
    \item Le fasi successive della pipeline (join, probe) eseguono anch'esse su GPU
\end{itemize}


\subsection{Stabilità delle misure}

\begin{table}[h]
\centering
\begin{tabular}{@{}lccc@{}}
\toprule
\textbf{Impl.} & \textbf{CoV medio} & \textbf{CoV max} & \textbf{Misure} \\
\midrule
Baseline & 1.98\% & 4.64\% & 20 \\
Autovec  & 0.39\% & 3.27\% & 15 \\
AVX2     & 3.46\% & 4.41\% & 20 \\
\bottomrule
\end{tabular}
\caption{Coefficiente di variazione $\text{CoV} = \sigma / \tilde{x}$ su 55 misure totali. Nessuna supera il 5\%.}
\end{table}

L'autovec è la più stabile perché GCC ottimizza lo scheduling delle istruzioni. Il CoV è maggiore a $N$ piccolo (4.64\% a $N=1$M per il baseline) dove la granularità del timer domina, e scende sotto 0.1\% per $N \geq 50$M.


% ############################################################################
\section{Correttezza e Validazione}
% ############################################################################

\subsection{Strategia a tre livelli}

\begin{enumerate}
    \item \textbf{Checksum FNV-1a} (per $N$ grandi): tutti i binari calcolano lo stesso hash dell'output. Probabilità di falso positivo $< 2^{-64}$.
    \item \textbf{Element-wise} (per $N \leq 32$): stampa completa e confronto manuale.
    \item \textbf{Verifica integrata} nel binario AVX2: esegue sia scalare che SIMD, confronta checksum, riporta prima posizione di mismatch.
\end{enumerate}

\subsection{Risultati di correttezza}
\begin{itemize}
    \item Checksum identici tra baseline, autovec, AVX2 e CUDA per \textbf{tutte} le 9 combinazioni $(N, P)$ testate
    \item Element-wise: verificato manualmente per $N=16, P=4$
    \item Distribuzione: $\max / \text{atteso} \leq 1.005$ per tutte le configurazioni
\end{itemize}


% ############################################################################
\section{Consigli, Pitfall e Scelte di Portabilità}
% ############################################################################

\begin{warnbox}[title=\textbf{Pitfall implementativi}]
\begin{enumerate}
    \item \textbf{\texttt{<filesystem>}}: non disponibile su tutte le toolchain. Usiamo API POSIX (\texttt{stat()}, \texttt{opendir()}).
    \item \textbf{\texttt{std::aligned\_alloc}}: non disponibile su macOS Apple libc++. Usiamo \texttt{posix\_memalign}.
    \item \textbf{\texttt{key\_t}}: conflitto con POSIX \texttt{key\_t}. Usiamo \texttt{spm\_key\_t}.
    \item \textbf{AVX2 su ARM}: il Makefile usa \texttt{uname -m} per abilitare \texttt{-mavx2} solo su x86\_64.
    \item \textbf{\texttt{make clean} in SLURM}: non fare \texttt{make clean} in script batch che girano in parallelo con altri job --- cancella i binari dell'altro job.
    \item \textbf{Allineamento}: \texttt{\_mm256\_load\_si256} richiede 32 byte. Usiamo \texttt{alloc\_aligned<T>(N, 32)}.
\end{enumerate}
\end{warnbox}

\begin{tipbox}[title=\textbf{Best practices}]
\begin{enumerate}
    \item Testa sempre con $N$ piccolo ($N=16$) prima di benchmark grandi.
    \item Usa \texttt{\_\_builtin\_ctz(P)} per calcolare $\log_2(P)$ (P potenza di 2).
    \item Warmup: 1 esecuzione prima delle misurazioni per popolare cache e TLB.
    \item \texttt{-march=native} su node09 abilita AVX2, FMA, e tutte le estensioni supportate.
    \item Compila con \texttt{-Wall -Wextra}: nessun warning nel nostro codice.
\end{enumerate}
\end{tipbox}


% ############################################################################
\section{Riferimenti}
% ############################################################################

\begin{enumerate}[label={[\arabic*]}]
    \item P.\ Ferragina, \textit{Pearls of Algorithm Engineering}, Cambridge University Press, 2023. Cap.\ 8: The Dictionary Problem, §8.2--8.3.
    \item D.\ Knuth, \textit{The Art of Computer Programming, Vol.\ 3: Sorting and Searching}, 2nd ed., 1998. §6.4 (Fibonacci hashing).
    \item Lezioni SPM 2025/26:
    \begin{itemize}
        \item \textbf{L5\&6}: Shared Memory --- von Neumann bottleneck (slide 3--5), modello Roofline (slide 47--51), operational intensity, ridge point
        \item \textbf{L7\&8}: SIMD on CPU --- speedup teorico (slide 7), intrinsics vs istruzioni (slide 8), SoA layout (slide 30, 39), condizioni auto-vectorization (slide 32--34, 40)
        \item \textbf{L9}: SIMT on GPU --- specs A30 (slide 24), memory coalescing (slide 28)
        \item \textbf{L10}: Metrics and Laws --- speedup/efficiency (slide 5--6), Amdahl's law (slide 24--28), overhead (slide 31--32)
        \item \textbf{L4}: SLURM --- gestione job, partizioni, allocazione esclusiva
    \end{itemize}
    \item M.\ Dietzfelbinger et al., ``A Reliable Randomized Algorithm for the Closest-Pair Problem'', \textit{J.\ Algorithms}, 1997.
    \item Intel, \textit{Intrinsics Guide}: \url{https://www.intel.com/content/www/us/en/docs/intrinsics-guide/index.html}
    \item S.\ Blackman, D.\ Vigna, ``Scrambled Linear Pseudorandom Number Generators'', \textit{ACM Trans.\ Math.\ Softw.}, 2021 (xoshiro256**).
\end{enumerate}


\end{document}
